\section{Einleitung}

\subsection{Projektüberblick}

TaskGrid+ ist ein verteiltes Aufgabenverwaltungssystem, das im Rahmen der Lehrveranstaltung
\textit{Verteilte Systeme} an der Technischen Hochschule entwickelt wurde.
Das System verteilt rechenintensive Aufgaben (Tasks) auf mehrere spezialisierte Worker-Prozesse
und gibt die Ergebnisse an anfragende Clients zurück.
Alle Komponenten kommunizieren ausschließlich über das gRPC-Protokoll;
Konfiguration erfolgt vollständig über Umgebungsvariablen.
Das Gesamtsystem wird mit Docker und Docker Compose containerisiert und betrieben.

\subsection{Systemziele}

Das System verfolgt folgende Kernziele:

\begin{itemize}
  \item \textbf{Entkopplung}: Dispatcher, Worker und Namensdienst kennen weder
        Adressen noch Implementierungsdetails der jeweils anderen Komponenten.
        Die Adressauflösung erfolgt ausschließlich dynamisch über den Namensdienst.
  \item \textbf{Skalierbarkeit}: Worker-Instanzen lassen sich zur Laufzeit
        starten und registrieren sich selbstständig beim Namensdienst.
        Pro Task-Typ können beliebig viele Instanzen parallel betrieben werden.
  \item \textbf{Robustheit}: Timeout-Erkennung und Retry-Mechanismus stellen sicher,
        dass ein ausgefallener Worker nicht zum dauerhaften Verlust eines Tasks führt.
  \item \textbf{Nachvollziehbarkeit}: Jede Zustandsänderung eines Tasks wird
        strukturiert geloggt mit \texttt{request\_id}, \texttt{task\_id} und Event-Name.
\end{itemize}

\subsection{Eigener Aufgabenbereich}

Die Implementierung wurde arbeitsteilig durchgeführt.
Diese Dokumentation beschreibt ausschließlich die selbst
bearbeiteten Issues:

\begin{table}[H]
\centering
\begin{tabular}{ll}
\toprule
\textbf{Issues} & \textbf{Komponente} \\
\midrule
\#13 -- \#21 & Dispatcher (Implementierung, Phase 2) \\
\#34 -- \#38 & Docker / Containerisierung (Phase 3) \\
\bottomrule
\end{tabular}
\caption{Bearbeitete Issues}
\end{table}

Weitere Komponenten (Namensdienst, Worker-Task-Handler, Client) wurden
von anderen Teammitgliedern implementiert und sind nicht Gegenstand
dieser Dokumentation.

\subsection{Aufbau der Dokumentation}

Die Dokumentation folgt dem Aufbau des Systems:
Abschnitt~\ref{sec:systemarchitektur} beschreibt die Gesamtarchitektur und
Komponentenstruktur.
Abschnitt~\ref{sec:kommunikation} erläutert das gRPC-Kommunikationsprotokoll.
Abschnitt~\ref{sec:datenmodell} dokumentiert das Task-Datenmodell und die Zustandsmaschine.
Abschnitt~\ref{sec:namensdienst} behandelt Service Discovery und Worker-Verwaltung.
Abschnitt~\ref{sec:technologien} beschreibt die Containerisierung mit Docker.
Abschnitt~\ref{sec:tests} belegt die Korrektheit durch Tests.
Abschnitt~\ref{sec:adr} dokumentiert die wichtigsten Architekturentscheidungen (ADRs).
